\section{Set-of-Mark：全图编号辅助定位}

\subsection{方法原理}

Set-of-Mark（SoM）是 Visual Prompting 的代表性工作，方法极其简单但非常有效：

\begin{enumerate}
    \item 对原始图片做全图实例分割（使用 SAM 的自动分割功能）
    \item 在每个分割区域的中心点打上数字编号
    \item 将带编号的新图片连同文本问题一起输入 VLM
    \item VLM 通过选择编号来回答定位问题
\end{enumerate}

\begin{importantbox}{SoM 的设计假设}
SoM 的有效性基于一个关键假设：VLM 在预训练阶段已经接触了大量带标注的 labelled figures、charts 和 textbooks。这些数据中天然包含编号、箭头等标记，因此 VLM 对这类视觉标记有很强的理解能力。SoM 利用这种能力，将图像分割为语义区域，赋予每个区域可读的 ID，让 VLM 能够"看图说话"。
\end{importantbox}

\subsection{优势分析}

SoM 的两大核心优势：

\begin{enumerate}
    \item \textbf{降低推理难度}：将直接生成坐标的回归问题变为选编号的分类问题
    \item \textbf{辅助注意力聚焦}：分割后的编号标记帮助模型聚焦到具体区域，稳定推理过程中的注意力
\end{enumerate}

\subsection{本章小结}

SoM 通过全图分割+编号标注的简单策略，将物体定位问题转化为选择题，显著降低 VLM 的推理难度，是 Visual Prompting 的经典范例。


